\documentclass[12pt,a4paper]{article}

\usepackage[utf8]{inputenc}
\usepackage[T1]{fontenc}
\usepackage[margin=2.5cm]{geometry}
\usepackage{hyperref}
\usepackage{microtype}
\usepackage{parskip}
\usepackage{lmodern}
\usepackage{amsmath}

\hypersetup{colorlinks=true, urlcolor=blue!60!black}

\pagestyle{empty}

\begin{document}

\begin{flushright}
Kundai Farai Sachikonye \\
Auf der Lichtung 19, 82178 Puchheim, Germany \\
\href{https://github.com/fullscreen-triangle}{github.com/fullscreen-triangle} \\
\texttt{kundai.sachikonye@bitspark.com} \\
(+49) 1626386344 \\[6pt]
20 May 2026
\end{flushright}

\vspace{1em}

PD Dr.\ med.\ Dr.\ rer.\ physiol.\ Peter Dietrich \\
Medizinische Klinik 1 \\
Universitätsklinikum Erlangen \\
Ref.\ he\_030\_PD

\vspace{1.5em}

\textbf{Re: Doktorand (m/w/d) --- Bioinformatik, Hepatologie / Onkologie (Ref.\ he\_030\_PD)}

\vspace{1em}

Dear PD Dr.\ Dietrich,

I am applying for the doctoral position (Ref.\ he\_030\_PD) in bioinformatic
analysis of human and murine tissues.
The combination of single-cell and bulk RNA sequencing analysis with a
hepatology and oncology focus describes precisely the analytical domain
for which I have developed and validated computational tools over the
past four years.

\textbf{RNA-seq analysis infrastructure.}
The \href{https://github.com/fullscreen-triangle/gospel}{Gospel} framework
handles the full RNA-seq analysis pipeline: read alignment and quantification,
differential expression analysis across conditions and cohorts, splicing analysis,
and integration with variant calling outputs for somatic mutation-expression
correlation.
Validated against 1000~Genomes Phase~3, gnomAD~v3.1.2, and UK~Biobank,
it processes $10^6$ variants per second with ClinVar, KEGG, and Reactome
pathway annotation built in.
For the comparison of human and murine transcriptomes --- a methodological
challenge specific to this project --- Gospel handles species-aware analysis
with ortholog mapping, differential gene expression across species, and
cross-species pathway conservation testing.
These comparisons are central to translational validation studies:
establishing which molecular phenotypes identified in murine models
are conserved in human patient tissue.

\textbf{Single-cell sequencing analysis.}
For single-cell RNA-seq, my primary contribution is the statistical
inference layer: variational Bayes posterior approximation at
$\mathcal{O}(10^6)$ latent variable scale, probabilistic graphical
models for multi-omics integration, and Gaussian mixture models for
cell population structure.
These methods sit directly below the standard Scanpy/Seurat pipeline ---
Leiden clustering and UMAP are dimensionality reduction steps whose
statistical behaviour I understand from the underlying geometry.
For hepatology-specific single-cell analysis this matters concretely:
the liver contains many rare cell types (hepatic stellate cells,
Kupffer cells, sinusoidal endothelial cells, periportal vs.\ pericentral
hepatocytes) whose identification requires probabilistic models that
handle low cell counts robustly, not just nearest-neighbour clustering.
RNA velocity and diffusion pseudotime --- trajectory inference for
fibrogenesis progression, tumour microenvironment remodelling, or
immune cell recruitment sequences --- are directly supported by the
Gaussian process infrastructure I have built.

\textbf{Hepatocellular carcinoma immunobiology.}
The oncology component of this project most likely involves the
hepatic tumour microenvironment: immune evasion, cytokine signalling
between tumour and stromal cells, or HCC neoantigen biology.
The \href{https://syndrome-seven.vercel.app/tools}{Syndrome} framework
I have developed computes spectral complementarity between molecular
sequences to predict MHC-I/II binding (AUC\,$0.81$ vs.\ NetMHCpan\,$0.68$
on 2,847 IEDB peptides), identify neoantigens from somatic mutation
data, and score immune escape variants via single amino acid substitution scanning.
For an RNA-seq study of HCC, this provides a direct bridge between
gene expression data and immunopeptidological function: identifying
which tumour-expressed transcripts generate immunogenic peptides,
which mutant peptides escape T cell recognition, and how the immune
state of the tumour microenvironment (quantified as the immune activation
dimension $D_\mathrm{Imm}$ of the disease state vector $\mathbf{D}\in[0,1]^8$)
correlates with treatment response.

\textbf{Molecular biology and in vivo experience.}
My computational background is primary, but it is grounded in wet lab
experience.
At Universitätsklinikum Hamburg-Eppendorf (2017--2018) I worked directly
in the Glycomics and Proteomics laboratory: automated LC-MS/MS and
MALDI-TOF pipeline operation, clinical sample processing, and data
quality control for multi-centre clinical cohort studies.
At TUM / Bayerisches Forschungsinstitut (2019--2021) I maintained
computational pipelines processing MS omics data from clinical cohorts,
working closely with the wet lab teams producing the data.
I understand sample preparation constraints, how technical variation
from different dissection or fixation protocols propagates into
sequencing data quality, and how to flag analytical artefacts that
originate in the biological protocol rather than the analysis.
The murine tissue component of this project is an opportunity to extend
that grounding to in vivo experimental design, which I pursue actively.

\textbf{Background.}
I hold an MSc in Computational Biology (Universität Tübingen, 2018--2019)
and an MSc in Pharmaceutical Biotechnology (HAW Hamburg, 2014--2017),
with doctoral research in Computational Lipidomics at TUM / Bayerisches
Forschungsinstitut (2019--2021).
I am legally resident in Germany with full work authorisation.
English is native; German is conversational.
I am available from 1 July 2026.

\vspace{1.5em}
Yours sincerely,

\vspace{2.5em}
\textbf{Kundai Farai Sachikonye}

\end{document}
